% preambule.tex — Paquets, mise en page, couleurs et commandes du mémoire.
% Compilation : pdflatex + biber (Overleaf : compilateur pdfLaTeX, biber détecté
% automatiquement ; en local : latexmk -pdf memoire.tex).

\usepackage[T1]{fontenc}
\usepackage[utf8]{inputenc}
\usepackage{lmodern}
\usepackage[french]{babel}
\usepackage{microtype}
\usepackage[a4paper, margin=2.5cm, headheight=14pt]{geometry}
\usepackage{setspace}
\setstretch{1.12}

\usepackage{amsmath, amssymb}
\usepackage[locale=FR]{siunitx}
\usepackage{graphicx}
\usepackage{xcolor}
\usepackage{booktabs, tabularx, array}
\usepackage{enumitem}
\usepackage{csquotes}
\usepackage[font=small, labelfont=bf]{caption}
\usepackage{subcaption}
\usepackage{float}
\usepackage{tikz}
\usetikzlibrary{babel, arrows.meta, positioning, calc, decorations.pathmorphing, shapes.geometric}
\usepackage{pgfplots}
\pgfplotsset{compat=1.18}
\usepackage{algorithm}
\usepackage{algpseudocode}
\usepackage[most]{tcolorbox}
\usepackage{titlesec}
\usepackage{fancyhdr}
\usepackage[backend=biber, style=authoryear-comp, maxcitenames=2, uniquelist=false]{biblatex}
\addbibresource{bibliographie.bib}
\usepackage[hidelinks]{hyperref}
\usepackage[nameinlink, noabbrev, french]{cleveref}

% Couleurs : reprises de l'application (accent bleu, convention des vecteurs)
\definecolor{accent}{HTML}{1976D2}
\definecolor{soft}{HTML}{E3F2FD}
\definecolor{poids}{HTML}{D32F2F}
\definecolor{normale}{HTML}{2E7D32}
\definecolor{tension}{HTML}{8E24AA}
\definecolor{corde}{HTML}{8D6E63}
\definecolor{eau}{HTML}{1E88E5}
\definecolor{metal}{HTML}{546E7A}
\definecolor{brouillon}{HTML}{FFF8E1}
\definecolor{brouillonbord}{HTML}{F9A825}

% Titres de chapitres sobres, numéro en couleur d'accent
\titleformat{\chapter}[display]
  {\normalfont\bfseries\color{accent}}
  {\large\MakeUppercase{\chaptertitlename}\ \thechapter}
  {6pt}
  {\Huge\color{black}}
  [\vspace{2pt}{\color{accent}\titlerule[1.2pt]}]
\titlespacing*{\chapter}{0pt}{-10pt}{24pt}
\titleformat{\section}{\normalfont\Large\bfseries}{\color{accent}\thesection}{0.8em}{}
\titleformat{\subsection}{\normalfont\large\bfseries}{\color{accent}\thesubsection}{0.8em}{}

\pagestyle{fancy}
\fancyhf{}
\fancyhead[L]{\small\nouppercase{\leftmark}}
\fancyhead[R]{\small\thepage}
\renewcommand{\headrulewidth}{0.4pt}
\fancypagestyle{plain}{\fancyhf{}\fancyfoot[C]{\small\thepage}\renewcommand{\headrulewidth}{0pt}}

\setlist{itemsep=0.25em, topsep=0.4em}

% Commandes
\newcommand{\fichier}[1]{\texttt{#1}}
\newcommand{\vect}[1]{\vec{#1}}

% Encadré « à rédiger » du squelette : objectif de la section, contenu attendu,
% nombre de pages visé. Supprimer l'encadré une fois la section rédigée.
\newtcolorbox{arediger}[1]{enhanced, breakable, colback=brouillon, colframe=brouillonbord,
  boxrule=0.6pt, arc=2pt, left=6pt, right=6pt, top=4pt, bottom=4pt,
  fonttitle=\small\bfseries, coltitle=black, title={À rédiger \hfill \(\approx\) #1}}

% Encadré de point clé dans le texte rédigé
\newtcolorbox{pointcle}{enhanced, colback=soft, colframe=accent, boxrule=0.6pt, arc=2pt,
  left=6pt, right=6pt, top=4pt, bottom=4pt}

% Algorithmes en français
\floatname{algorithm}{Algorithme}
\renewcommand{\listalgorithmname}{Liste des algorithmes}
\algrenewcommand\algorithmicrequire{\textbf{Entrée :}}
\algrenewcommand\algorithmicensure{\textbf{Sortie :}}
\algrenewcommand\algorithmicfor{\textbf{pour}}
\algrenewcommand\algorithmicdo{\textbf{faire}}
\algrenewcommand\algorithmicend{\textbf{fin}}
\algrenewcommand\algorithmicif{\textbf{si}}
\algrenewcommand\algorithmicthen{\textbf{alors}}
\algrenewcommand\algorithmicwhile{\textbf{tant que}}
